% Options for packages loaded elsewhere
\PassOptionsToPackage{unicode}{hyperref}
\PassOptionsToPackage{hyphens}{url}
%
\documentclass[
  ignorenonframetext,
]{beamer}
\usepackage{pgfpages}
\setbeamertemplate{caption}[numbered]
\setbeamertemplate{caption label separator}{: }
\setbeamercolor{caption name}{fg=normal text.fg}
\beamertemplatenavigationsymbolsempty
% Prevent slide breaks in the middle of a paragraph
\widowpenalties 1 10000
\raggedbottom
\setbeamertemplate{part page}{
  \centering
  \begin{beamercolorbox}[sep=16pt,center]{part title}
    \usebeamerfont{part title}\insertpart\par
  \end{beamercolorbox}
}
\setbeamertemplate{section page}{
  \centering
  \begin{beamercolorbox}[sep=12pt,center]{part title}
    \usebeamerfont{section title}\insertsection\par
  \end{beamercolorbox}
}
\setbeamertemplate{subsection page}{
  \centering
  \begin{beamercolorbox}[sep=8pt,center]{part title}
    \usebeamerfont{subsection title}\insertsubsection\par
  \end{beamercolorbox}
}
\AtBeginPart{
  \frame{\partpage}
}
\AtBeginSection{
  \ifbibliography
  \else
    \frame{\sectionpage}
  \fi
}
\AtBeginSubsection{
  \frame{\subsectionpage}
}
\usepackage{amsmath,amssymb}
\usepackage{lmodern}
\usepackage{iftex}
\ifPDFTeX
  \usepackage[T1]{fontenc}
  \usepackage[utf8]{inputenc}
  \usepackage{textcomp} % provide euro and other symbols
\else % if luatex or xetex
  \usepackage{unicode-math}
  \defaultfontfeatures{Scale=MatchLowercase}
  \defaultfontfeatures[\rmfamily]{Ligatures=TeX,Scale=1}
\fi
\usetheme[]{Montpellier}
% Use upquote if available, for straight quotes in verbatim environments
\IfFileExists{upquote.sty}{\usepackage{upquote}}{}
\IfFileExists{microtype.sty}{% use microtype if available
  \usepackage[]{microtype}
  \UseMicrotypeSet[protrusion]{basicmath} % disable protrusion for tt fonts
}{}
\makeatletter
\@ifundefined{KOMAClassName}{% if non-KOMA class
  \IfFileExists{parskip.sty}{%
    \usepackage{parskip}
  }{% else
    \setlength{\parindent}{0pt}
    \setlength{\parskip}{6pt plus 2pt minus 1pt}}
}{% if KOMA class
  \KOMAoptions{parskip=half}}
\makeatother
\usepackage{xcolor}
\IfFileExists{xurl.sty}{\usepackage{xurl}}{} % add URL line breaks if available
\IfFileExists{bookmark.sty}{\usepackage{bookmark}}{\usepackage{hyperref}}
\hypersetup{
  hidelinks,
  pdfcreator={LaTeX via pandoc}}
\urlstyle{same} % disable monospaced font for URLs
\newif\ifbibliography
\usepackage{color}
\usepackage{fancyvrb}
\newcommand{\VerbBar}{|}
\newcommand{\VERB}{\Verb[commandchars=\\\{\}]}
\DefineVerbatimEnvironment{Highlighting}{Verbatim}{commandchars=\\\{\}}
% Add ',fontsize=\small' for more characters per line
\usepackage{framed}
\definecolor{shadecolor}{RGB}{248,248,248}
\newenvironment{Shaded}{\begin{snugshade}}{\end{snugshade}}
\newcommand{\AlertTok}[1]{\textcolor[rgb]{0.94,0.16,0.16}{#1}}
\newcommand{\AnnotationTok}[1]{\textcolor[rgb]{0.56,0.35,0.01}{\textbf{\textit{#1}}}}
\newcommand{\AttributeTok}[1]{\textcolor[rgb]{0.77,0.63,0.00}{#1}}
\newcommand{\BaseNTok}[1]{\textcolor[rgb]{0.00,0.00,0.81}{#1}}
\newcommand{\BuiltInTok}[1]{#1}
\newcommand{\CharTok}[1]{\textcolor[rgb]{0.31,0.60,0.02}{#1}}
\newcommand{\CommentTok}[1]{\textcolor[rgb]{0.56,0.35,0.01}{\textit{#1}}}
\newcommand{\CommentVarTok}[1]{\textcolor[rgb]{0.56,0.35,0.01}{\textbf{\textit{#1}}}}
\newcommand{\ConstantTok}[1]{\textcolor[rgb]{0.00,0.00,0.00}{#1}}
\newcommand{\ControlFlowTok}[1]{\textcolor[rgb]{0.13,0.29,0.53}{\textbf{#1}}}
\newcommand{\DataTypeTok}[1]{\textcolor[rgb]{0.13,0.29,0.53}{#1}}
\newcommand{\DecValTok}[1]{\textcolor[rgb]{0.00,0.00,0.81}{#1}}
\newcommand{\DocumentationTok}[1]{\textcolor[rgb]{0.56,0.35,0.01}{\textbf{\textit{#1}}}}
\newcommand{\ErrorTok}[1]{\textcolor[rgb]{0.64,0.00,0.00}{\textbf{#1}}}
\newcommand{\ExtensionTok}[1]{#1}
\newcommand{\FloatTok}[1]{\textcolor[rgb]{0.00,0.00,0.81}{#1}}
\newcommand{\FunctionTok}[1]{\textcolor[rgb]{0.00,0.00,0.00}{#1}}
\newcommand{\ImportTok}[1]{#1}
\newcommand{\InformationTok}[1]{\textcolor[rgb]{0.56,0.35,0.01}{\textbf{\textit{#1}}}}
\newcommand{\KeywordTok}[1]{\textcolor[rgb]{0.13,0.29,0.53}{\textbf{#1}}}
\newcommand{\NormalTok}[1]{#1}
\newcommand{\OperatorTok}[1]{\textcolor[rgb]{0.81,0.36,0.00}{\textbf{#1}}}
\newcommand{\OtherTok}[1]{\textcolor[rgb]{0.56,0.35,0.01}{#1}}
\newcommand{\PreprocessorTok}[1]{\textcolor[rgb]{0.56,0.35,0.01}{\textit{#1}}}
\newcommand{\RegionMarkerTok}[1]{#1}
\newcommand{\SpecialCharTok}[1]{\textcolor[rgb]{0.00,0.00,0.00}{#1}}
\newcommand{\SpecialStringTok}[1]{\textcolor[rgb]{0.31,0.60,0.02}{#1}}
\newcommand{\StringTok}[1]{\textcolor[rgb]{0.31,0.60,0.02}{#1}}
\newcommand{\VariableTok}[1]{\textcolor[rgb]{0.00,0.00,0.00}{#1}}
\newcommand{\VerbatimStringTok}[1]{\textcolor[rgb]{0.31,0.60,0.02}{#1}}
\newcommand{\WarningTok}[1]{\textcolor[rgb]{0.56,0.35,0.01}{\textbf{\textit{#1}}}}
\setlength{\emergencystretch}{3em} % prevent overfull lines
\providecommand{\tightlist}{%
  \setlength{\itemsep}{0pt}\setlength{\parskip}{0pt}}
\setcounter{secnumdepth}{-\maxdimen} % remove section numbering
%\documentclass[unicode,10pt]{beamer}% 'unicode'が必要
\usepackage{luatexja}% 日本語を使う場合は必須
%\usepackage[japanese]{babel}	
%\usefonttheme[onlymath]{serif}
%\usepackage[T1]{fontenc} %これを使うと、ギリシャ文字が出ない
%\usepackage{textcomp}
%\usepackage[scale=1.0]{tgheros} % Sans serif font
%\usepackage[scaled]{beramono} % Monospace font
%\usepackage{luatexja-otf}
%\renewcommand{\kanjifamilydefault}{\gtdefault}
%\usepackage[haranoaji]{luatexja-preset}

% スライド番号の表示 %%%%%%%%%%%%%%%%%%%
\makeatletter
\setbeamertemplate{footline}{%
  \leavevmode%
  \hbox{%
  \begin{beamercolorbox}[wd=.5\paperwidth,ht=2.25ex,dp=1ex,right]{date in head/foot}%
    \insertframenumber{} \hspace*{2ex} % new version without total frames
  \end{beamercolorbox}}%
  \vskip0pt%
}
\makeatother
% スライド番号の表示 %%%%%%%%%%%%%%%%%%%
\ifLuaTeX
  \usepackage{selnolig}  % disable illegal ligatures
\fi

\title{第2章: 因果関係\\
(2.2. Rでデータを部分集合化する)}
\subtitle{今井耕介 著\\
『社会科学のためのデータ分析入門 (QSS)』}
\author{}
\date{\vspace{-2.5em}2026-03-09}

\begin{document}
\frame{\titlepage}

\hypertarget{rux3067ux30c7ux30fcux30bfux3092ux90e8ux5206ux96c6ux5408ux5316ux3059ux308b}{%
\section{2.2
Rでデータを部分集合化する}\label{rux3067ux30c7ux30fcux30bfux3092ux90e8ux5206ux96c6ux5408ux5316ux3059ux308b}}

\begin{frame}[fragile]{データの抽出（部分集合化）の目的}
\protect\hypertarget{ux30c7ux30fcux30bfux306eux62bdux51faux90e8ux5206ux96c6ux5408ux5316ux306eux76eeux7684}{}
\begin{itemize}
\tightlist
\item
  大規模なデータセット全体を分析するだけでなく、特定の条件を満たす「部分集合」に注目したい場合が多い。

  \begin{itemize}
  \tightlist
  \item
    例：男性だけのデータ、特定の州のデータ、ある年齢層のデータなど。
  \end{itemize}
\item
  Rでは、\textbf{論理演算子}と\textbf{インデックス（\texttt{{[}{]}} や
  \texttt{subset()}）}を組み合わせて、柔軟にデータを抽出できる。
\end{itemize}
\end{frame}

\hypertarget{ux8ad6ux7406ux6f14ux7b97ux5b50-logical-operators}{%
\section{2.2.1 論理演算子 (Logical
Operators)}\label{ux8ad6ux7406ux6f14ux7b97ux5b50-logical-operators}}

\begin{frame}[fragile]{基本的な論理演算子}
\protect\hypertarget{ux57faux672cux7684ux306aux8ad6ux7406ux6f14ux7b97ux5b50}{}
\begin{itemize}
\tightlist
\item
  2つの値を比較し、\texttt{TRUE}（真）か \texttt{FALSE}（偽）を返す。
\end{itemize}

\scriptsize

\begin{Shaded}
\begin{Highlighting}[]
\CommentTok{\# 1. 等しい (==)}
\DecValTok{5} \SpecialCharTok{==} \DecValTok{5}
\end{Highlighting}
\end{Shaded}

\begin{verbatim}
## [1] TRUE
\end{verbatim}

\begin{Shaded}
\begin{Highlighting}[]
\DecValTok{5} \SpecialCharTok{==} \DecValTok{3}
\end{Highlighting}
\end{Shaded}

\begin{verbatim}
## [1] FALSE
\end{verbatim}

\begin{Shaded}
\begin{Highlighting}[]
\CommentTok{\# 2. 等しくない (!=)}
\DecValTok{5} \SpecialCharTok{!=} \DecValTok{3}
\end{Highlighting}
\end{Shaded}

\begin{verbatim}
## [1] TRUE
\end{verbatim}

\begin{Shaded}
\begin{Highlighting}[]
\CommentTok{\# 3. 大なり、小なり (\textgreater{}, \textless{}, \textgreater{}=, \textless{}=)}
\DecValTok{5} \SpecialCharTok{\textgreater{}} \DecValTok{3}
\end{Highlighting}
\end{Shaded}

\begin{verbatim}
## [1] TRUE
\end{verbatim}

\normalsize
\end{frame}

\begin{frame}[fragile]{複数の条件の組み合わせ}
\protect\hypertarget{ux8907ux6570ux306eux6761ux4ef6ux306eux7d44ux307fux5408ux308fux305b}{}
\begin{itemize}
\tightlist
\item
  「かつ（AND）」や「または（OR）」を使って条件を繋げる。
\end{itemize}

\scriptsize

\begin{Shaded}
\begin{Highlighting}[]
\CommentTok{\# 1. かつ (\&): 両方の条件が真なら TRUE}
\NormalTok{(}\DecValTok{5} \SpecialCharTok{\textgreater{}} \DecValTok{3}\NormalTok{) }\SpecialCharTok{\&}\NormalTok{ (}\DecValTok{2} \SpecialCharTok{\textless{}} \DecValTok{4}\NormalTok{)}
\end{Highlighting}
\end{Shaded}

\begin{verbatim}
## [1] TRUE
\end{verbatim}

\begin{Shaded}
\begin{Highlighting}[]
\CommentTok{\# 2. または (|): どちらか一方が真なら TRUE}
\NormalTok{(}\DecValTok{5} \SpecialCharTok{\textgreater{}} \DecValTok{3}\NormalTok{) }\SpecialCharTok{|}\NormalTok{ (}\DecValTok{2} \SpecialCharTok{\textgreater{}} \DecValTok{4}\NormalTok{)}
\end{Highlighting}
\end{Shaded}

\begin{verbatim}
## [1] TRUE
\end{verbatim}

\normalsize
\end{frame}

\hypertarget{ux30d9ux30afux30c8ux30ebux306bux5bfeux3059ux308bux8ad6ux7406ux6f14ux7b97}{%
\section{2.2.2
ベクトルに対する論理演算}\label{ux30d9ux30afux30c8ux30ebux306bux5bfeux3059ux308bux8ad6ux7406ux6f14ux7b97}}

\begin{frame}[fragile]{ベクトルの要素ごとに比較する}
\protect\hypertarget{ux30d9ux30afux30c8ux30ebux306eux8981ux7d20ux3054ux3068ux306bux6bd4ux8f03ux3059ux308b}{}
\begin{itemize}
\tightlist
\item
  Rの論理演算はベクトル全体に対して一度に行われ、各要素についての
  \texttt{TRUE/FALSE} が返される。
\end{itemize}

\scriptsize

\begin{Shaded}
\begin{Highlighting}[]
\CommentTok{\# サンプルの読み込み}
\NormalTok{resume }\OtherTok{\textless{}{-}} \FunctionTok{read.csv}\NormalTok{(}\StringTok{"resume.csv"}\NormalTok{)}

\CommentTok{\# race 変数の最初の5つの要素が "black" かどうかを確認}
\FunctionTok{head}\NormalTok{(resume}\SpecialCharTok{$}\NormalTok{race }\SpecialCharTok{==} \StringTok{"black"}\NormalTok{, }\AttributeTok{n =} \DecValTok{5}\NormalTok{)}
\end{Highlighting}
\end{Shaded}

\begin{verbatim}
## [1] FALSE FALSE  TRUE  TRUE FALSE
\end{verbatim}

\begin{Shaded}
\begin{Highlighting}[]
\CommentTok{\# コールバック (call) が 1 であった行の TRUE/FALSE}
\FunctionTok{head}\NormalTok{(resume}\SpecialCharTok{$}\NormalTok{call }\SpecialCharTok{==} \DecValTok{1}\NormalTok{, }\AttributeTok{n =} \DecValTok{5}\NormalTok{)}
\end{Highlighting}
\end{Shaded}

\begin{verbatim}
## [1] FALSE FALSE FALSE FALSE FALSE
\end{verbatim}

\normalsize
\end{frame}

\hypertarget{ux30c7ux30fcux30bfux306eux62bdux51fa-indexing}{%
\section{2.2.3 データの抽出
(Indexing)}\label{ux30c7ux30fcux30bfux306eux62bdux51fa-indexing}}

\begin{frame}[fragile]{角括弧 \texttt{{[}{]}} を使った抽出 (1)}
\protect\hypertarget{ux89d2ux62ecux5f27-ux3092ux4f7fux3063ux305fux62bdux51fa-1}{}
\begin{itemize}
\tightlist
\item
  \texttt{データ{[}行の条件,\ 列の指定{]}}
  という形式でデータを抽出する。
\end{itemize}

\scriptsize

\begin{Shaded}
\begin{Highlighting}[]
\CommentTok{\# 1. 黒人 (race == "black") の行だけを抽出}
\CommentTok{\# カンマの後は空にすると「全ての列」を意味する}
\NormalTok{resumeB }\OtherTok{\textless{}{-}}\NormalTok{ resume[resume}\SpecialCharTok{$}\NormalTok{race }\SpecialCharTok{==} \StringTok{"black"}\NormalTok{, ]}

\CommentTok{\# 2. 白人 (race == "white") の行だけを抽出}
\NormalTok{resumeW }\OtherTok{\textless{}{-}}\NormalTok{ resume[resume}\SpecialCharTok{$}\NormalTok{race }\SpecialCharTok{==} \StringTok{"white"}\NormalTok{, ]}

\CommentTok{\# 抽出したデータの行数を確認}
\FunctionTok{nrow}\NormalTok{(resumeB)}
\end{Highlighting}
\end{Shaded}

\begin{verbatim}
## [1] 2435
\end{verbatim}

\begin{Shaded}
\begin{Highlighting}[]
\FunctionTok{nrow}\NormalTok{(resumeW)}
\end{Highlighting}
\end{Shaded}

\begin{verbatim}
## [1] 2435
\end{verbatim}

\normalsize
\end{frame}

\begin{frame}[fragile]{角括弧 \texttt{{[}{]}} を使った抽出 (2)}
\protect\hypertarget{ux89d2ux62ecux5f27-ux3092ux4f7fux3063ux305fux62bdux51fa-2}{}
\begin{itemize}
\tightlist
\item
  抽出したサブセットに対して、統計量を計算する。
\end{itemize}

\scriptsize

\begin{Shaded}
\begin{Highlighting}[]
\CommentTok{\# 黒人のコールバック率（平均）}
\FunctionTok{mean}\NormalTok{(resumeB}\SpecialCharTok{$}\NormalTok{call)}
\end{Highlighting}
\end{Shaded}

\begin{verbatim}
## [1] 0.06447639
\end{verbatim}

\begin{Shaded}
\begin{Highlighting}[]
\CommentTok{\# 白人のコールバック率（平均）}
\FunctionTok{mean}\NormalTok{(resumeW}\SpecialCharTok{$}\NormalTok{call)}
\end{Highlighting}
\end{Shaded}

\begin{verbatim}
## [1] 0.09650924
\end{verbatim}

\normalsize
\end{frame}

\hypertarget{subset-ux95a2ux6570ux306eux6d3bux7528}{%
\section{2.2.4 subset()
関数の活用}\label{subset-ux95a2ux6570ux306eux6d3bux7528}}

\begin{frame}[fragile]{subset() による直感的な抽出}
\protect\hypertarget{subset-ux306bux3088ux308bux76f4ux611fux7684ux306aux62bdux51fa}{}
\begin{itemize}
\tightlist
\item
  \texttt{subset(データ,\ subset\ =\ 条件)}
  という書き方で、より読みやすく抽出できる。
\end{itemize}

\scriptsize

\begin{Shaded}
\begin{Highlighting}[]
\CommentTok{\# 黒人かつ男性 (sex == "male") のデータを抽出}
\NormalTok{resumeBm }\OtherTok{\textless{}{-}} \FunctionTok{subset}\NormalTok{(resume, }\AttributeTok{subset =}\NormalTok{ (race }\SpecialCharTok{==} \StringTok{"black"} \SpecialCharTok{\&}\NormalTok{ sex }\SpecialCharTok{==} \StringTok{"male"}\NormalTok{))}

\CommentTok{\# データの最初の数行を表示}
\FunctionTok{head}\NormalTok{(resumeBm, }\AttributeTok{n =} \DecValTok{3}\NormalTok{)}
\end{Highlighting}
\end{Shaded}

\begin{verbatim}
##    firstname  sex  race call
## 10    Tyrone male black    0
## 21     Leroy male black    0
## 42    Tyrone male black    0
\end{verbatim}

\begin{Shaded}
\begin{Highlighting}[]
\CommentTok{\# 黒人男性のコールバック率}
\FunctionTok{mean}\NormalTok{(resumeBm}\SpecialCharTok{$}\NormalTok{call)}
\end{Highlighting}
\end{Shaded}

\begin{verbatim}
## [1] 0.0582878
\end{verbatim}

\normalsize
\end{frame}

\hypertarget{ux307eux3068ux3081}{%
\section{2.2.5 まとめ}\label{ux307eux3068ux3081}}

\begin{frame}[fragile]{このセクションのまとめ}
\protect\hypertarget{ux3053ux306eux30bbux30afux30b7ux30e7ux30f3ux306eux307eux3068ux3081}{}
\begin{itemize}
\tightlist
\item
  \textbf{論理演算子}: \texttt{==}, \texttt{!=},
  \texttt{\textgreater{}}, \texttt{\textless{}}, \texttt{\&},
  \texttt{\textbar{}} を使って、データが条件を満たすかを判定する。
\item
  \textbf{データの抽出}:

  \begin{itemize}
  \tightlist
  \item
    \textbf{角括弧 \texttt{{[}{]}}}:
    \texttt{resume{[}resume\$race\ ==\ "black",\ {]}} のように記述。
  \item
    \textbf{subset() 関数}: \texttt{subset(resume,\ race\ ==\ "black")}
    のように、よりシンプルに記述可能。
  \end{itemize}
\item
  \textbf{因果推論への応用}:
  グループごとにデータを切り分けることで、処置群と対照群の平均を簡単に比較できる。
\end{itemize}
\end{frame}

\end{document}
